% Generated draft: requires mathematical and source review.
\documentclass[11pt]{article}
\usepackage[margin=25mm]{geometry}
\usepackage{fontspec}
\setmainfont{Latin Modern Roman}
\usepackage{amsmath,amssymb,mathtools,unicode-math}
\setmathfont{Latin Modern Math}
\usepackage{enumitem,needspace,xurl,hyperref}
\setlength{\parindent}{0pt}
\setlength{\parskip}{5pt}
\setlength{\emergencystretch}{3em}
\newcommand{\N}{\mathbb{N}}
\newcommand{\Z}{\mathbb{Z}}
\newcommand{\Q}{\mathbb{Q}}
\newcommand{\R}{\mathbb{R}}
\newcommand{\C}{\mathbb{C}}
\newcommand{\F}{\mathbb{F}}
\newcommand{\A}{\mathbb{A}}
\newcommand{\PP}{\mathbb{P}}
\newcommand{\E}{\mathbb{E}}
\newcommand{\eps}{\varepsilon}
\newcommand{\dd}{\,\mathrm{d}}
\begin{document}
\section*{UnsolvedMath ID: 20000461}
\textbf{Source problem code:} catalogue entry
\subsection{Definitions and mathematical statement}
Let \(E/\Q\) have a nonidentity point \(P\in E(\Q)[3]\) and \(\operatorname{rank}E(\Q)=0\). Let \(K/\Q\) be cyclic of degree three, of conductor \(f\), and let \(\nu(f)\) count its distinct prime divisors. The conventional finite \(3\)-Selmer group is
\[
 \operatorname{Sel}_3(E/K)=\ker\left(H^1(K,E[3])\longrightarrow
 \prod_vH^1(K_v,E)\right),
\]
where the maps impose the local Kummer conditions. Chantal David's question is whether, for every such \(E,K\),
\[
 3^{\nu(f)-1}\mid\#\operatorname{Sel}_3(E/K),
 \quad\text{equivalently}\quad
 \dim_{\F_3}\operatorname{Sel}_3(E/K)\ge\nu(f)-1.
\]
The cardinality sign is part of the supplied verification's repair; divisibility is not applied to a group without taking its order.

The supplied partial theorem retains rational \(3\)-torsion and the cyclic cubic extension but does not require rank zero. Put \(C=\langle P\rangle\), \(E'=E/C\), and \(\phi:E\to E'\). The \(\phi\)-Selmer group is the subgroup of \(H^1(K,C)\) satisfying the local Kummer conditions from \(E'(K_v)/\phi E(K_v)\). Identify \(C\) with \(\Z/3\) using \(P\). The everywhere-unramified character space is
\[
 U_K=\operatorname{Hom}(\operatorname{Cl}_K,\F_3)\subseteq H^1(K,C),\qquad
 r_3(\operatorname{Cl}_K)=\dim_{\F_3}(\operatorname{Cl}_K/3\operatorname{Cl}_K).
\]
Define the localization defect
\[
 \delta(E,K)=r_3(\operatorname{Cl}_K)
 -\dim_{\F_3}(U_K\cap\operatorname{Sel}_\phi(E/K)).
\]
Let \(N_E\) be the conductor of \(E\), let \(\omega(m)\) count distinct prime divisors of a positive integer \(m\), and let \(\rho(E,K)\) be the dimension of the span in \(\operatorname{Cl}_K/3\operatorname{Cl}_K\) of the prime-ideal classes above rational primes dividing \(3N_E\). The report asserts
\[
 0\le\delta(E,K)\le\rho(E,K)\le2\omega(3N_E),
\]
that \(\operatorname{Sel}_\phi(E/K)\) injects into \(\operatorname{Sel}_3(E/K)\), using the dual kernel \(\mu_3\) and the fact that cyclic cubic fields are totally real, and that genus theory gives \(r_3(\operatorname{Cl}_K)\ge\nu(f)-1\). Consequently its reported bound is
\[
 \dim_{\F_3}\operatorname{Sel}_3(E/K)
 \ge\nu(f)-1-\delta(E,K)
 \ge\nu(f)-1-2\omega(3N_E).
\]
A negative right side is merely a vacuous dimension lower bound, not divisibility by a negative power interpreted as an integer. The exact requested bound follows under \(\delta(E,K)=0\), in particular if every prime class above \(3N_E\) belongs to \(3\operatorname{Cl}_K\). This conditional statement does not assert vanishing of the defect for every \(K\).

The historical remarks motivate the question by cyclic-cubic genus-theory class-number divisibility \(3^{d-1}\mid h_K\), with \(d\) the number of ramified discriminant primes, and report progress of the form \(3^{d-c_E}\) with \(c_E\) depending on \(E\). They explicitly question whether that progress concerns an algebraic \(L\)-value, a \(3\)-primary Selmer group or a \(3\)-descent group, and mention computations of special \(L\)-values. Those computations are not silently treated as proofs for the finite Selmer group defined above. Poonen proposes analogues for other primes or base number fields, without a precise replacement statement.
\subsection{Short English statement}
Does every rank-zero elliptic curve over \(\Q\) with rational three-torsion acquire a \(3\)-Selmer group of order divisible by \(3^{\nu(f)-1}\) over a cyclic cubic field? The supplied report gives a weaker conductor-dependent lower bound and recovers the exact bound when a specified localization defect vanishes.
\subsection{Sources}
\begin{itemize}
\item Local UnsolvedMath registry, permanent ID \(20000461\), source code \texttt{catalogue entry}; AIM algebraic number theory corpus, \texttt{aim-algebraic-number-theory-notes.json}, notes[112]. Attribution: Ulam AI / UnsolvedMath Contributors, CC BY 4.0.
\item AIM \emph{Rational and integral points on higher dimensional varieties}, Question 28 attributed to Chantal David, with remarks including Poonen: \url{https://aimath.org/WWN/qptsurface2/qptsurface2.pdf}. The supplied verification reports that official HTML/PDF restore the lost cardinality sign; the \(P\ne O\) and exponent typography are likewise made explicit. The remarks end with ``Problem/''.
\item Original literature field is empty; no full reference or proof is supplied for the historical progress or the later local-defect theorem.
\item Registry dated August 14, 2026: statement exact, status \texttt{partially\_solved}, classification \texttt{PARTIAL-PROGRESS}; attempt 1 records problem status open, partial result, rigor proved, novelty confidence low. The defect decomposition and bounds are machine-generated research aids, not peer reviewed, requiring independent verification before citation.
\end{itemize}
\end{document}
